\section{Adaptive sampling parameters and precision termination}\label{at:details}
The primitive theta/Mellin identity with root number one yields
\[
 \Lambda_q(s)=\sum_{n\ge1}\chi_q(n)n^a
 \int_1^\infty e^{-\pi n^2t/q}
 \bigl(t^{(s+a)/2-1}+t^{(1-s+a)/2-1}\bigr)\,dt.
\]
On $D=[-1,2]+i[-4/3,4/3]$, both exponents $p$ satisfy
$|\Re p|\le2$ and $|\Im p|\le2/3$.
For exact $x\ge0$, put $E(x)=2^{-\lfloor x\rfloor}$.
The elementary inequality $e>2$ gives
$e^{-x}\le e^{-\lfloor x\rfloor}\le E(x)$ and
$E(x)\le2e^{-x\log2}$.
Use this solely as a pointwise nonincreasing upper bound; it is not
multiplicative and has no asserted exponential integral identity.

Integration by parts gives
$\int_1^\infty t^2e^{-xt}\,dt=e^{-x}Q(x)$ for
$Q(x)=x^{-1}+2x^{-2}+2x^{-3}$.
For $a=0$ the bound $n^ae^{-\delta n^2/2}\le C_0=1$ is immediate.
For $a=1$, $n\le(1+n^2)/2$ and
$n^2e^{-\delta n^2/2}\le2/(e\delta)\le1/\delta$ give
$C_1=(1+1/\delta)/2$.
For $n=n_0+j$, the inequality
$n^2\ge n_0^2+2n_0j$ therefore bounds the omitted-$n$ sum by
\[
 2C_aQ(\delta n_0^2)
 \frac{e^{-\delta n_0^2/2}}{1-e^{-\delta n_0}}
 \le2C_aQ(\delta n_0^2)
 \frac{E(\delta n_0^2/2)}{1-E(\delta n_0)}=B_n.
\]
Take $n_0=N+1\ge2$; then $\delta n_0\ge6/5$ and the latter denominator
is positive. Substituting $t=Tu$ in each retained-$n$ time tail gives
\[
 B_t=\sum_{n=1}^N2n^aT^3E(\delta n^2T)Q(\delta n^2T).
\]
The two tails are disjoint because the omitted-$n$ tail uses the whole
integral and the time tail uses only retained $n$.

For $j=0,\ldots,J-1$ and $\ell=0,\ldots,3$ take
\[
 m=2^j\left(1+\frac{2\ell+1}{8}\right),\qquad h=\frac{2^j}{8},
 \qquad T=2^J.
\]
These panels cover $[1,T]$ and have $h/m\le1/9$.
On $|z-m|=\rho=m/2$, principal log is analytic,
$\Re z\ge m/2$, and $|\arg z|\le1$.
For $f(z)=e^{-\pi n^2z/q}z^p$, the angular multiplier is at most
$e^{2/3}\le3$, whence
\[
 |f(z)|\le3B(m)E(\delta n^2m/2),\qquad
 B(m)=\max\{(3m/2)^2,(m/2)^{-2},1\}.
\]
Cauchy's coefficient estimate, cancellation of odd terms, and geometric
summation yield the $B_Q$ in Theorem~\ref{at:sample}.
Including zero-character or later discarded terms in that sum is harmless.
Gaussian, time, and geometric decay prove termination of the sequential
first-success searches with $B_n,B_t,B_Q\le\eta/8$.
For fixed q and $L=\max\{1,\log(1/\eta)\}$ their asymptotic scales can be
chosen as $N=O(\sqrt L)$, $T=O(L)$, $J=O(\log L)$, and $K=O(L)$.
Polynomial prefactors in $N,T,J$ contribute only $O(\log L)$ to the
logarithm of the error estimates; this is a truncation-size statement,
not a rational-arithmetic bit-complexity bound.

For real $t\in[m-h,m+h]$, $t\ge1$ and $t^{\Re p}\le(m+h)^2$.
The bound for a whole pair, including both powers, is therefore
$4hn^a(m+h)^2E(\delta n^2(m-h))$.
There are $P=4JN$ potential pairs; summing only bounds at most
$\eta/(8P)$ gives $B_{\rm discard}\le\eta/8$.
No partial recurrence or isolated exponential is discarded.
For active pairs retain the coefficient recurrence
\[
 r_{-1}=0,\quad r_0=1,\quad
 r_{j+1}=\frac{(h/m)(p-\lambda m-j)r_j-
                   \lambda m(h/m)^2r_{j-1}}{j+1},
 \qquad\lambda=\pi n^2/q.
\]
Coefficient comparison in $t f'=(-\lambda t+p)f$ proves that the
integrated Taylor polynomial is
\[
 2he^{p\log m-\lambda m}\sum_{j=0}^K\frac{r_{2j}}{2j+1}.
\]
Outward arithmetic encloses its exact value, including the signed
character weight and both powers.

Use an isolated p-bit dyadic context, with all primitive remainders
refined as p increases. Put $\delta_p=2^{-p-16}$.
Each arctangent series separately stops when its next-term bound is at
most $\delta_p$, giving total weighted analytic tail width at most
$20\delta_p$ in Machin's formula, before outward dyadic rounding.
Each reduced logarithm series separately has tail at most $\delta_p$;
for $x=2^kv$, the combined bound is $(1+|k|)\delta_p$, before rounding.
For the complex exponential, after reduction to $A\le1/2$, choose the
Taylor order so the bound
\[
 \frac{A^{H+1}}{(H+1)!}\frac1{1-A/(H+2)}.
\]
is at most $\delta_p$. Choose H by finite exact search; its denominator
is positive. Add a symmetric exponential error enclosure to each coordinate
before outward squaring. The fixed guard16 is not itself an acceptance
bound: the final complete rectangle radius must satisfy
$R_{\rm arith}\le\eta/4$. All interval caches must include p and preserve
exact input-type guards. The successful p and retry precisions/radii are
recorded. Individual series lengths are determined by these, the exact
inputs and the pinned source and are reconstructible by replay; they are
not explicitly stored.

For fixed N,J,K and exact target point, the expression has finitely many
operations and only nonzero fixed integer denominators. Each primitive
input is bounded for sufficiently large p; exponential halving counts
are consequently bounded. Exact primitive tails and dyadic rounding
widths tend to zero. Induction through the finite continuous expression,
including finitely many outward squarings, gives $R_{\rm arith}\to0$.
Doubling p after a failed comparison terminates. The argument also covers
true roots and nodes regenerated at every p, since their widths tend to zero.
A fixed nondegenerate s rectangle has no corresponding shrinking-input
claim.

For serialization, initialize
\[
 B=\max\{32,\operatorname{bitlen}(\operatorname{den}\eta)
             -\operatorname{bitlen}(\operatorname{num}\eta)+8\}
\]
and increase B until $5\,2^{-B}\le\eta/8$.
Write $\widehat x=2^{-B}\lceil2^Bx\rceil$ for each of the five raw
components. Their exact reported sum is at most
$3\eta/4+5\,2^{-B}\le7\eta/8$, and every component remains an upper
bound. The output records B, the raw interval radius and complete rectangle,
N/J/K, discard threshold/pair counts, exact target metadata and source hashes.
These reconstruct all raw analytic bounds and pair sets without serializing
huge exact tail denominators. For the practical requests below B is87
and114, respectively. Exact endpoints remain rational; the mathematical
specification assumes an integer representation sufficient for the requested
size. Dynamic rounding removes the old fixed reporting floor. For rational s, validate its exact coordinates in closed D before
conversion: an outward dyadic enclosure of $\Im s=4/3$ extends beyond D,
but the analytic bounds concern the true point. Circle nodes satisfy the
strict exact bounds $-7/10\le\Re s\le17/10$, $|\Im s|\le6/5$.

The exact parameter-design checker gives
\[
\begin{array}{c|c|rrrr|r}
\eta&q&N&J&T&K&\text{active pairs}\\\hline
10^{-24}&3&12&7&128&19&71\\
10^{-24}&4&14&8&256&20&65\\
10^{-24}&5&16&8&256&20&111\\
10^{-32}&3&14&7&128&25&84\\
10^{-32}&4&16&8&256&26&72\\
10^{-32}&5&18&8&256&26&129
\end{array}
\]
with exact rational budgets and no special-function oracle.
The arithmetic precision is selected from the actual interval radius,
not inferred from this table. This extension retains the existing
classical theta identity; it adds no zero-location, positive determinant,
all-conductor, Lean or GRH claim. A downstream adaptive readout requires
its own node-count, alias, degree, Fourier and optimization choices.
